\documentclass[12pt,a4paper]{article}

% Package imports for Serbian Cyrillic support
\usepackage[T2A]{fontenc}
\usepackage[utf8]{inputenc}
\usepackage[serbian]{babel}

% Use Times font for better Cyrillic support (usually available)
\usepackage{times}

% Other packages
\usepackage{geometry}
\usepackage{hyperref}
\usepackage{fancyhdr}
\usepackage{titlesec}
\usepackage{enumitem}
\usepackage{listings}
\usepackage{xcolor}
\usepackage{graphicx}
\usepackage{tikz}
\usepackage{booktabs}
\usepackage{amsmath}
\usepackage{amssymb}

% Page geometry
\geometry{
    left=2.5cm,
    right=2.5cm,
    top=3cm,
    bottom=3cm
}

% Hyperlink setup
\hypersetup{
    colorlinks=true,
    linkcolor=blue,
    filecolor=magenta,      
    urlcolor=cyan,
    pdftitle={Спецификација пројекта 2025/26},
    pdfpagemode=FullScreen,
}

% Header and footer
\pagestyle{fancy}
\fancyhf{}
\fancyhead[L]{Спецификација пројекта 2025/26}
\fancyfoot[C]{\thepage}

% Custom commands for special formatting
\definecolor{emerald}{RGB}{0,128,0}         % classic emerald
\definecolor{purple}{RGB}{128,0,128}     % classic purple

\newcommand{\processitem}[1]{\textbf{#1}}
\newcommand{\optional}[1]{\textcolor{blue}{#1}}
\newcommand{\bonus}[1]{\textcolor{emerald}{#1}}
\newcommand{\extra}[1]{\textcolor{purple}{#1}}

% Code listing setup
\lstset{
    basicstyle=\ttfamily\small,
    breaklines=true,
    frame=single,
    backgroundcolor=\color{gray!10}
}

% Document information
\title{\textbf{Спецификација пројекта 2025/26}}
\author{Upravljanje poslovnim procesima}
\date{\vspace{-\baselineskip}}

\begin{document}

\maketitle

\section{Увод}

Задатак пројекта је развити process-dirven апликацију која подржава пословни процес кодификације закона.

Ради се о систему за кодификацију закона у коме сарађују Централна агенција за кодификацију и више различитих влада. Пословни процеси су дефинисани тако да нагласе колаборацију између организација.

У реалним околностима ови процеси су компликованији са строжијим роковима. За потребе пројекта смо поједноставили неке од корака да би акценат био на самим пословним процесима.

\subsection{Претпоставке за окружење}

\begin{itemize}
    \item Мора постојати иницијалини регистар влада (минимум 5).
    \item Корисници који могу да користе систем морају бити унапред регистровани.
    \item Основни подаци регистрованих корисника су: име и презиме корисника, имеил, корисничко име, шифра, адреса, место, поштански број.
    \item Владе морају имати регистрованог корисника за приступ систему (адвокати који наступају у њихово име).
\end{itemize}

\section{Процеси које систем подржава}

У овом поглављу описани су главни процеси које апликација треба да подржи.

\subsection{Процес кодификације закона (колаборативни процес)}

Процес описује ток од тренутка када Иницијална влада пошаље захтјев за кодификацију до објављивања нове верзије закона или прекида поступка. Подразумијева се да за дату Владу већ постоји активна конфигурација услуге (правила прегледа, процјене утицаја и рокова - описано у другом процесу).

\subsubsection{Увид у статистичке податке агенције}

\extra{Адвокати који су менаџери у агенцији имају увид у историјске статистичке податке агенције. На основу ових података агенција гледа да унаприједи своје услуге и убрза свој процес рада. Статистички преглед укључује:}

\begin{itemize}
    \item \extra{Увид у све тренутне захтјеве за кодификацију, у ком су стању, за коју владу и који је рок за израду. Захтјеви су сортирани по крајњем року за завршетак кодификације.}
    \item \extra{Информације о: просјечном трајању кодификације свих до сад успјешно одрађених кодификација, просјечно трајање кодификације уколико је било финалног прегледа од стране владе, просјечно трајање уколико је била процјена утицаја, просјечно трајање уколико је било и финални преглед и процјена утицаја.}
    \item \extra{Графикон који приказује све успјешно завршене кодификације са временом трајања.}
    \item \extra{Графикон који приказује однос успјешних и неуспјешних кодификација.}
    \item \extra{Графикон који приказује кодификације по тренутном сатусу.}
    \item \extra{Графикон који приказује број кодификација по домену закона.}
    \item \extra{Графикон који приказује за сваку кодификацију колико времена је провела у агенцији, а колико се чекало на одговоре од влада.}
\end{itemize}

\subsubsection{Иницирање захтјева}

Процес иницира запослени у Иницијалној влади уносом основних података о закону (назив, да ли је у питању нови закон или измјене и допуне, домени на које се односи, ниво управе), степена хитности (нормалан или хитан) и ознаке да ли закон може да утиче на друге владе. Уз захтјев се отпрема један или више докумената (нпр. MS Word документ(и), прилози, опциони скенирани потписани документ у PDF формату). Захтјев се као порука шаље Агенцији.

\subsubsection{Иницијална провјера у Агенцији}

Захтјев преузима запослени у Агенцији (примарни адвокат или службеник за пријем) и провјерава да ли је документација комплетна и да ли је захтјев у опсегу услуга које Агенција пружа. Ако захтјев није исправан или није у опсегу, Агенција га одбија, шаље Влади образложење и процес се завршава стањем „одбијен од стране агенције". Ако је захтјев прихваћен, додјељује се примарни адвокат и процес прелази у фазу појашњења и кодификације. \optional{Адвокат који се додјељује је онај који тренутно не ради ни једну кодификацију, уколико нема таквог адвоката, бира се адвокат чији је процес у стању чекања и најближи крају.}

\subsubsection{Фаза појашњења}

Примарни адвокат прегледа све достављене документе/цјелине. За сваки дио за који су потребна појашњења адвокат формулише питања за Владу. Питања се групишу по документу или логичкој цјелини. Агенција затим шаље Влади поруку са прикупљеним питањима. Запослени адвокати у Влади припремају одговоре по документу/цјелини (сваки адвокат у влади може да погледа и да одговор), а само овлашћени запослени може да пошаље коначне одговоре назад Агенцији (запослени који је иницијално започео процес). Примарни адвокат прегледа одговоре; ако су појашњења довољна, прелази се на кодификацију. Ако нису, формирају се нова питања и циклус (слање питања – одговори – провјера) се понавља. Уколико ни након 3 итерације нису добијена потребна појашњења, Агенција може да одбије захтјев.

\optional{Запослени у влади који учествују у фази појашњења бирају се на основу домена за који су специјализовани. Тј. могу да учествују сви адвокати владе који су специјализовани за домене за које се кодификују закони.}

\subsubsection{Кодификација}

Примарни адвокат врши кодификацију: уграђује нове чланове у постојећи текст, мијења или брише постојеће одредбе, ажурира упућивања на друге законе и слично. Након завршетка обраде свих дијелова, формира се јединствена кодификована верзија закона.

\optional{За потребе пројекта симулираћемо кодификацију тако што ће систем да споји све PDF документе који се кодификују у један јединствен документ.}

\subsubsection{Унутрашњи преглед у Агенцији}

Кодификована верзија се шаље на преглед другом адвокату (адвокат за преглед).

\optional{Систем предлаже адвоката за преглед. Адвокат се бира или на основу домена (ако је раније кодификовао законе из предложеног домена), ако нема таквих адвоката онда на основу владе (ако је већ тој влади радио кодификацију) и на крају на основу нивоа управе (ако је већ кодификовао за другу владу на истом нивоу управе). Уколико нема такав адвокат, бира се насумично адвокат који тренутно не ради ни једну кодификацију. Уколико сви адвокати раде на некој кодификацији тренутно, бира се насумично један чији је процес у стању чекања и најближи крају.}

Адвокат за преглед може да:
\begin{enumerate}[label=\alph*)]
    \item одобри верзију, у ком случају процес прелази на одлуку о процјени утицаја;
    \item врати коментаре примарном адвокату.
\end{enumerate}

Примарни адвокат уноси исправке и шаље нову верзију на преглед. Циклус: кодификација – преглед – исправке, се понавља док адвокат за преглед не одобри текст. Исти адвокат прегледа текст сваки наредни пут, одабир адвоката за преглед се ради само пред прву итерацију.

\subsubsection{Одлука о процјени утицаја}

Након унутрашњег прегледа, Агенција провјерава да ли је потребна процјена утицаја на друге владе, на основу ознаке да закон има међувладин утицај, домена закона и конфигурације услуге за дату Владу. Ако процјена утицаја није потребна, процес наставља ка прегледу од стране Иницијалне владе (тачка 7) (ако је предвиђен) или директно ка објављивању (тачка 9). Процес се аутоматски покреће из процеса кодификације када су испуњени услови за процјену утицаја (означено да закон има међувладин утицај, релевантни домени, конфигурација услуге).

\paragraph{Покретање процеса.}

Агенција припрема пакет за процјену утицаја који садржи сажетак кодификованог закона, списак чланова или одредби које могу утицати на погођене владе, опис могућих посљедица и рок за одговор. Један такав пакет може да се користи за све погођене владе, са могућношћу додавања специфичних напомена по влади. Агенција бира које су владе погођене из свих доступних влада у систему.

\bonus{За 2 додатна бода, аутоматизовати корак одабира погођених влада. На основу улазних параметара, препоручити кориснику иницијални списак погођених влада које онда корисник може да доради.}

\paragraph{Слање захтјева погођеним владама.}

Агенција затим шаље свакој погођеној влади поруку са захтјевом за процјену утицаја.

\extra{Погођене владе на својој страни виде преглед свих активних захтјева са основним статусом (примљен, у прегледу, одговор послат), ко је рецензент, која влада је издала закон, за које домене.}

\paragraph{Преглед у погођеним владама.}

У свакој погођеној влади захтјев се додјељује рецензентима утицаја. Рецензент(и) анализирају пакет и утврђују прихватљивост утицаја, потребу за измјенама или постојање блокирајућих проблема. Ако учествује више рецензената, њихови закључци се консолидују у јединствен званичан одговор те владе.

\optional{Систем бира рецензенте аутоматски на основу домена. Сваки адвокат у влади има листу домена у којима је стручан. На основу ових домена, одабрати рецензенте аутоматски.}

Званичан одговор може бити:
\begin{itemize}
    \item одобрено без примједби (ако нема ни једне примједбе),
    \item одобрено са примједбама (више од 50\% рецензената је одобрило и постоји бар једна примједба),
    \item неодобрено (блокирајуће примједбе) (испод 50\% рецензената је одобрило).
\end{itemize}
Одговор се шаље Агенцији као порука.

\paragraph{Рокови за одговор.}

За сваки захтјев постоји рок до кога погођена влада треба да одговори. Ако рок истекне без одговора, систем шаље подсјетник погођеној влади и обавјештава Агенцију. Уколико ни након подсјетника одговор не стигне, Агенција може да тај исход третира као посебну категорију („без одговора") и да га на одговарајући начин укључи у консолидацију резултата. Рок за одговор је 2 дана пред истек званичног рока за завршетак кодификације.

\paragraph{Консолидација у Агенцији.}

Како одговори пристижу, Агенција за сваку владу биљежи исход и коментаре. Када су сви одговори примљени или су истекли рокови (укључујући сценарије без одговора), Агенција консолидује резултате у један укупни исход за цијели процес процјене утицаја. Могући консолидовани исходи су:
\begin{itemize}
    \item успјешно, без блокирајућих примједби (редован наставак кодификације),
    \item успјешно, са препорукама за измјене (враћа се на процес измјене, без потребе да се поново ради процјена утицаја),
    \item неуспјешно (постоје блокирајуће примједбе) (враћа се на процес измејене, са поновном провјером утицаја. нове измјене се шаљу само владама које нису одобриле кодификацију),
    \item није завршено (недостају одговори или није постигнут јасан исход). (конфигурација подешава како се односи у овом сценарију, ако није специфицирано, наставља се процес).
\end{itemize}
На основу тога Агенција припрема сажет извјештај.

\subsubsection{Преглед од стране Иницијалне владе}

Уколико конфигурација услуге предвиђа финални преглед, Агенција шаље кодификовану верзију закона Иницијалној влади на одобрење. У овом прегледу влада има увид у све релевантне податке, новокодификован закон, као и извјештај о погођеним владама (уколико постоји), Влада може да:
\begin{enumerate}[label=\alph*)]
    \item одобри верзију, након чега процес наставља ка објављивању;
    \item пошаље примједбе и затражи измјене.
\end{enumerate}

У случају примједби, примарни адвокат врши измјене, по потреби поново пролази унутрашњи преглед (и поновну процјену утицаја), па затим шаље нову верзију назад Влади. Циклус прегледа и измјена се понавља док Влада не одобри верзију или док се процес не прекине. Преглед кодификованог закона у влади врше сви адвокати који су учествовали у фази појашњења.

\optional{Уколико није било фазе појашњења, преглед врше адвокати из владе који су специјализовани за домене за које су кодификовани закони.}

\optional{Агенција генерише извјештај о кодификацији и шаље га влади. Извјештај је ПДФ документ који саджи све релевантне податке о кодификацији, сам текст новокодификованог закона, као и извјештај о погођеним владама (уколико постоји).}

\subsubsection{Повлачење захтјева од стране Владе}

Иницијална влада у одређеним фазама (нпр. током појашњења или прегледа) може да повуче захтјев за кодификацију из политичких или других разлога. У том случају Влада шаље Агенцији поруку о повлачењу, сви активни задаци повезани са тим захтјевом се завршавају, а процес завршава у стању „повучен од стране владе".

\subsubsection{Објављивање кодификованог закона}

Када су завршене све потребне фазе (поjашњења, кодификација, унутрашњи преглед, евентуална процјена утицаја и преглед Владе), кодификована верзија је спремна за објављивање. Запослени у Агенцији задужен за објаву врши завршну техничку провјеру, и позива спољашњи сервис за објављивање или индексирање (није потребно имплементирати у оквиру пројекта). Након успјешне објаве, Агенција обавјештава Владу и процес завршава у стању „успјешно објављен".

\subsubsection{Рокови и контролисан прекид}

Агенција има рок од 30 дана да заврши кодификацију укључујући све прегледа и процјене утицаја од стране владе.

\bonus{За додатних 3 бода: Агенција подразумјевано има рок од 5 радних дана да заврши кодификацију. На кључним мјестима процеса (чекање на појашњења, одговоре на преглед, резултат процјене утицаја) рок се паузира. У случају прекида због пробијања рокова, процес завршава у стању „прекинут због пробијања рокова".}

\subsection{Процес конфигурације услуге између Агенције и Владе}

Процес дефинише правила сарадње између Централне агенције за кодификацију и појединачне Владе. Резултат је активна конфигурација услуге која се касније користи у процесима кодификације и процјене утицаја. По потреби проширити могуће ставке за конфигурацију уколико има потребе за тим (нпр. није специфициран подразумјевана вриједност у самом процесу кодификације).

\subsubsection{Увид у статистичке податке агенције}

\extra{Влада има увид у историјске статистичке податке агенције. На основу ових података може да направи смислен приједлог конфигурације услуга. Статистички преглед укључује:}

\begin{itemize}
    \item \extra{Колико просјечно траје кодификација на основу свих до сада извршених кодификација свих влада које су на истом нивоу управе.}
    \item \extra{Колико просјечно траје кодификација уколико се тражи финални преглед, колико просјечно траје уколико се тражи процјена утицаја, колико просјечно траје ако се раде и финални прегелд и процјена утицаја.}
    \item \extra{Колико времена просјечно агенција чека на одговор за финални преглед.}
\end{itemize}
\extra{Све податке приказивати само на основу претходних кодификација за владе које су на истом нивоу управе.}

\subsubsection{1. Покретање захтјева од стране Владе}

Запослени у Влади покреће процес и попуњава форму за конфигурацију. Уносе се основни подаци о Влади (назив, ниво управе, контакт), типови прописа које Влада планира да шаље на кодификацију, опције прегледа (да ли увијек жели финални преглед прије објављивања, да ли жели увид у резултате процјене утицаја), подешавања процјене утицаја (за које домене је обавезна) и предложени рокови (рок за нормалну и хитну кодификацију, рок за одговор Владе на питања Агенције), итд. Након слања форме, захтјев се као порука шаље Агенцији.

\subsubsection{Иницијална обрада у Агенцији}

Захтјев преузима запослени Агенције и провјерава да ли су подаци комплетни, да ли су тражени рокови и опције реални и да ли је опсег у складу са услугама које Агенција пружа. Агенција може да затражи допуну или да припреми измијењени приједлог конфигурације.

\subsubsection{Предлог конфигурације од стране Агенције}

На основу захтјева Владе и интерних правила, Агенција припрема предлог конфигурације који се као порука шаље Влади, уз кратко образложење евентуалних измјена у односу на иницијални захтјев.

\subsubsection{Одговор Владе}

Надлежни у Влади прегледа предлог Агенције и може да:
\begin{enumerate}[label=\alph*)]
    \item прихвати предлог, након чега се прелази на активацију конфигурације;
    \item одбије или измјени предлог и формира нову конфигурацију (нпр. другачији рокови, додатни домени за процјену утицаја, итд.), која се шаље Агенцији као порука.
\end{enumerate}

Кораци 3 и 4 чине преговарачки циклус који се може више пута понављати: Агенција шаље нови предлог, Влада га одобрава или враћа са измјенама. Ако у том року од 30 дана или 3 итерације не дође до сагласности, процес се завршава као неуспјешан. Агенција обавјештава Владу да конфигурација није договорена и да услуга кодификације не може да се пружа под траженим условима. Процес завршава стањем „конфигурација није договорена". Влада без активне конфигурације не може да покреће редовне процесе кодификације закона.

\subsubsection{Активација конфигурације}

Када се Влада и Агенција усагласе да је приједлог реалан, креира се активна конфигурација услуге за ту Владу. Конфигурација садржи основне податке о начину извршавања процеса кодификације. И Влада и Агенција добијају обавјештење о успјешној активацији, а процес завршава стањем „успјешно активирана конфигурација".

\section{Задатак пројекта:}

\begin{itemize}
    \item Направити модел пословних података који омогућавају имплементацију процеса који подржава све горе наведене операције.
    \item Направити модел пословног процеса који омогућава извршење свих неопходних активности.
    \item Развити модел учесника у процесу и извршити одговарајуће доделе задатака.
    \item Водити рачуна да модел буде што флексибилнији, тј. да се додела појединачних таскова ради на основу података прикупљених током извршавања процеса. Ово посебно важи за доделу задатака адвокатима у владама и запосленим у агенцији.
    \item На местима где је то неопходно обезбедити повезивање са екстерним сервисима.
    \item Водити рачуна да сви процеси морају бити временски ограничени.
    \item Одређене активности обавља више корисника, ту употребити мулти-инстанце таскове.
    \item По потреби моделе процеса изделити на мање процесе како би се повећала прегледност и квалитет модела. Бодовање модела процеса ће зависити од квалитета и читљивости модела.
    \item Осигурати се да апликација неће отказати тако што ће се искористити одговарајући одговори на отказе.
    \item Обезбедити одговарајуће форме које су неопходне да се обаве задаци који се појављаују у разним фазама процеса.
    \item Развијени модел деплоyовати на Бонита engine.
    \item \optional{Обавјештења, напомене, слање порука и сличне акције из спецификације реализовати као слање имејла одговарајућим корисницима.}
    \item \extra{Реализовати ово као посебне апликације – направити application deployment descriptore. Један за владе и један за агенцију. Сваки треба да буде персонализован институцији у којој ће се користити са специфичним приказима који ће им дати увид у тренутно стање процеса и олакшати извођење.}
\end{itemize}

\section{Ограничења и бодовање:}

\begin{itemize}
    \item Пројекат се ради у тиму од два члана. Користи се \href{https://www.bonitasoft.com/old-versions}{Бонита 2023}.
    \item Пројекат укупно носи 45 бодова.
    \item За домаћи задатак је потребно направити модел података и моделовати организацију. Домаћи задатак носи 5 бодова.
    \item \textbf{Прва контролна тачка (14 бодова)} обухвата задатке 2.1.2 - 2.1.5 из процеса кодификације, као и задатке 2.2.2 - 2.2.5 из процеса конфигурације. Потребно је направити барем два процеса. Обезбедити покретање и извршавање ових процеса независно од улаза од којих они зависе (започети процесе са mock подацима за тестирање ако је то потребно). За прву контролну тачку сваки од процеса носи по 7 бодова. Не мора да радите мулти-инстанце задатке за ову конторлну тачку. Сваки задатак може бити додјељен једном кориснику, па ћете за другу конторлну тачку то измјенити да одговара спецификацији. За ову конторлну тачку процес кодификације не мора да зависи од активне конфигурације, можете користити неке подразумјеване вриједности.
    \item \textbf{Друга контролна тачка (16 бодова)} обухвата остале задатке из оба процеса. Исправити поједностављене задатке из прве контролне тачке. Увести мулти-инстанце задатке. Потребно је да се увежу процеси тако да рокови и начин извршавања процеса кодификације зависи од активне конфигурације коју влада тренутно има. Није потребно радити са конкторима за генерисање докумената и слање имејлова, то припада трећој контролној тачки.
    \item \textbf{Трећа контролна тачка (10 бодова)} подразумијева задатке означене \optional{плавом бојом}. Међу задацима означеним плавом бојом постоји 5 јединствених задатака од којих сваки носи по 2 бода. Неки од задатака се раде на више мјеста, нпр слање имејла, те је потребно одрадити задатак на свим мјестима за максимум бодова.
    \item Свака контролна тачка ће бити благовремено најављена и путем обавештења на Canvas-у и по потреби са додатним информацијама о изради.
\end{itemize}

Могуће је надокнадити изгубљене бодове. Надокнада је једино могућа израдом тачака 2.1.7 и/или 2.1.11 (рок од 5 дана са паузама) из процеса кодификације. То су дијелови текста означени \bonus{зеленом бојом}. Овим путем можете надокнадити 5 изгубљених бодова. Бодове можете надокнадити приликом одбране КТ3.

Укоилко не желите да одбраните пројекат путем контролних тачка, сви задаци који су опциони или додатни постају обавезни (нпр. рок од 5 дана са паузама). Додатно, потребно је урадити задатке означене \extra{љубичастом бојом} (нпр. направити application descriptor-е). Овај додатни рок за одбрану пројекта ће накнадно бити објављен и оквирно ће бити пред априлски рок.

\end{document}
